\documentclass[11pt,a4paper]{article}
\usepackage{lmodern}
\usepackage[T1]{fontenc}
\usepackage[utf8]{inputenc}
\usepackage[italian]{babel}
\usepackage[a4paper,margin=2.6cm,headheight=14pt]{geometry}
\usepackage{graphicx}
\usepackage{booktabs}
\usepackage{xcolor}
\usepackage{fancyhdr}
\usepackage{caption}
\usepackage{listings}
\usepackage{microtype}
\usepackage{longtable}
\usepackage[hidelinks]{hyperref}

\definecolor{accent}{HTML}{1F4E8C}
\definecolor{codebg}{HTML}{F4F6F8}
\definecolor{rule}{HTML}{C8D0D8}

\hypersetup{colorlinks=true,linkcolor=accent,urlcolor=accent,citecolor=accent}

\pagestyle{fancy}
\fancyhf{}
\fancyhead[L]{\footnotesize\textsc{Struttura del repository}}
\fancyhead[R]{\footnotesize\thepage}
\renewcommand{\headrulewidth}{0.4pt}

\captionsetup{font=small,labelfont={bf,color=accent},width=0.92\textwidth}

\lstset{
  basicstyle=\ttfamily\footnotesize,
  backgroundcolor=\color{codebg},
  frame=leftline, framerule=1.6pt, rulecolor=\color{accent},
  framexleftmargin=6pt, xleftmargin=8pt,
  breaklines=true, breakatwhitespace=false,
  columns=fullflexible, keepspaces=true,
  showstringspaces=false,
  aboveskip=9pt, belowskip=9pt,
  literate={à}{{\`a}}1 {è}{{\`e}}1 {é}{{\'e}}1 {ì}{{\`i}}1 {ò}{{\`o}}1 {ù}{{\`u}}1
           {À}{{\`A}}1 {È}{{\`E}}1 {°}{{\textdegree}}1 {–}{{-}}1 {—}{{---}}1
}

\newcommand{\us}{\,$\mu$s}
\setlength{\parskip}{0.45em}
\setlength{\parindent}{0pt}
\title{\vspace{-1.6cm}\textbf{Struttura del repository}\\[.2em]
\large Valutazione di OpenTelemetry per workload mixed-criticality}
\author{}
\date{\today}

\begin{document}
\maketitle
\vspace{-1.4cm}

\section*{Panoramica}
\addcontentsline{toc}{section}{Panoramica}

La radice del progetto \`e organizzata seguendo il flusso di lavoro: dal codice sorgente agli
eseguibili, da questi agli strumenti che li configurano ed eseguono, fino ai dati prodotti e ai
documenti che li interpretano. Le cartelle numerate rispettano l'ordine in cui il lavoro \`e stato
svolto.

\begin{center}\small
\begin{tabular}{p{3.1cm}p{1.3cm}p{9.9cm}}
\toprule
\textbf{cartella} & \textbf{peso} & \textbf{ruolo} \\
\midrule
\texttt{rt-app/} & 20\,MB & Il generatore di carico real-time strumentato, con i sorgenti e l'infrastruttura di compilazione. \\
\texttt{otel-installdir/} & 12\,MB & La libreria OpenTelemetry compilata e installata localmente, contro cui il generatore viene collegato. \\
\texttt{bin/} & 66\,MB & Gli eseguibili gi\`a compilati, uno per ciascuna configurazione di strumentazione. \\
\texttt{scripts/} & 188\,KB & Gli strumenti sviluppati per il progetto: preparazione della piattaforma, esecuzione delle campagne, analisi. \\
\texttt{1-configs/} & 24\,KB & Le configurazioni di carico di riferimento, in forma leggibile. \\
\texttt{2-DoE/} & 204\,MB & Tutti i dati sperimentali grezzi e i risultati dell'analisi. \\
\texttt{0-explore/} & 2.2\,MB & Le prove preliminari condotte a mano prima di automatizzare le misure. \\
\texttt{3-report/} & 776\,KB & I documenti di sintesi e i relativi sorgenti. \\
\bottomrule
\end{tabular}
\end{center}

\vspace{-.4em}
\section{\texttt{rt-app/} --- il generatore di carico}

Contiene il programma che genera il workload periodico e ne registra i tempi. \`E la base di
partenza fornita per il progetto, gi\`a strumentata con OpenTelemetry.

\begin{center}\small
\begin{tabular}{p{3.4cm}p{10.9cm}}
\toprule
\texttt{src/} & I sorgenti, dodici file fra implementazione e intestazioni. I moduli principali
riguardano il corpo del programma e l'esecuzione dei task, l'interpretazione del file di
configurazione, la gestione degli argomenti da riga di comando, la scrittura dei log e i gruppi di
task. Un'intestazione raccoglie le definizioni che governano la strumentazione. \\
\addlinespace
\texttt{libdl/} & Una piccola libreria di supporto che espone le chiamate di sistema per lo
scheduling con scadenza, non disponibili direttamente nella libreria standard. \\
\addlinespace
\texttt{doc/} & Materiale accessorio distribuito con il programma originale. \\
\addlinespace
file di build & \texttt{configure.ac}, \texttt{Makefile.am} e \texttt{autogen.sh} sono i sorgenti
dell'infrastruttura di compilazione; \texttt{configure}, i \texttt{Makefile} e i file di supporto
sono generati automaticamente e non vanno modificati a mano. \\
\bottomrule
\end{tabular}
\end{center}

La compilazione avviene in tre passi (\texttt{./autogen.sh}, \texttt{./configure},
\texttt{make}) e produce l'eseguibile in \texttt{src/}. I parametri che governano la strumentazione
vengono passati al compilatore in questa fase, non a esecuzione avviata: \`e la ragione per cui
esiste la cartella descritta di seguito.

\section{\texttt{bin/} --- gli eseguibili per ciascuna configurazione}

Poich\'e il comportamento della telemetria \`e fissato al momento della compilazione, ogni
configurazione sperimentale richiede un eseguibile distinto. Ricompilare a ogni esecuzione
sarebbe stato lento e, soprattutto, avrebbe reso difficile garantire che due esecuzioni della stessa
cella usassero davvero lo stesso binario. Gli eseguibili vengono quindi costruiti una sola volta e
conservati qui, indicizzati dalla combinazione di parametri che rappresentano.

Il nome codifica per intero la configurazione:

\begin{center}\ttfamily\large
rtapp\_t\underline{3}\_p\underline{1}\_s\underline{0}\_r\underline{0.0}\_e\underline{0}
\end{center}

\begin{center}\small
\begin{tabular}{clp{8.4cm}}
\toprule
\textbf{sigla} & \textbf{parametro} & \textbf{valori} \\
\midrule
\texttt{t} & livello di tracciamento & 0 nessuno, 1 esecuzione e thread, 2 anche le fasi, 3 anche ogni iterazione \\
\texttt{p} & modalit\`a di consegna & 0 differita (thread dedicato), 1 immediata (nel thread stesso) \\
\texttt{s} & criterio di campionamento & 0 sempre attivo, 1 proporzionale, 2 mai attivo \\
\texttt{r} & proporzione & valore usato quando il criterio \`e proporzionale \\
\texttt{e} & destinazione & 0 servizio di raccolta esterno, 1 uscita standard \\
\bottomrule
\end{tabular}
\end{center}

La cartella contiene attualmente quattordici eseguibili: undici corrispondenti alle celle delle tre
campagne, e tre prodotti durante le prove preliminari, riconoscibili dalla diversa convenzione di
denominazione.

\section{\texttt{scripts/} --- gli strumenti sviluppati per il progetto}

\`E la parte realizzata da noi. Si divide in due gruppi, secondo la funzione.

\subsection*{\texttt{scripts/utils\_isolation/} --- preparazione della piattaforma}

\begin{center}\small
\begin{tabular}{p{4.3cm}p{10.0cm}}
\toprule
\texttt{pin\_cpu\_freq.sh} & Fissa la frequenza del processore e disabilita la sovra-frequenza
automatica, agendo sui registri di controllo. Accetta i comandi \texttt{status}, \texttt{fix} e
\texttt{reset}. \\
\addlinespace
\texttt{isolate\_cpus.sh} & Crea la partizione fra CPU di servizio e CPU riservate agli esperimenti,
sposta su queste prime tutti i processi esistenti e riassegna le affinit\`a degli interrupt.
Salva lo stato precedente per poterlo ripristinare. \\
\addlinespace
\texttt{reset\_isolation.sh} & Annulla l'isolamento e ripristina le affinit\`a salvate,
riportando il sistema allo stato iniziale. \\
\bottomrule
\end{tabular}
\end{center}

Accanto ai primi due si trovano file con estensione \texttt{.orig}: conservano la versione di
partenza degli script, che abbiamo poi corretto ed esteso. Sono mantenuti per rendere tracciabili
le modifiche.

\subsection*{\texttt{scripts/measurements/} --- esecuzione e analisi}

\begin{center}\small
\begin{tabular}{p{4.3cm}p{10.0cm}}
\toprule
\texttt{gen\_config.py} & Genera la configurazione del carico a partire dai parametri (numero di
task di disturbo, durata, CPU, costante di calibrazione), garantendo che le celle differiscano solo
per ci\`o che si intende variare. Avverte se i task venissero collocati su CPU logiche dello stesso
core fisico. \\
\addlinespace
\texttt{test.sh} & Esegue una singola prova all'interno dell'area isolata, in una cartella dedicata
in cui salva anche la configurazione effettivamente usata. \\
\addlinespace
\texttt{run\_doe.sh} & Orchestra una campagna completa: costruisce gli eseguibili necessari, genera
le configurazioni, esegue le ripetizioni in ordine interlacciato, registra le condizioni ambientali
e verifica la validit\`a di ciascuna esecuzione. Accetta come argomento la campagna da eseguire. \\
\addlinespace
\texttt{analyze\_doe.py} & Percorre tutte le esecuzioni, ne estrae le variabili di risposta e le
unisce ai livelli dei fattori, producendo le tabelle dei risultati. \\
\addlinespace
\texttt{analyze\_block1/2/3.py} & Analisi specifiche di ciascuna campagna, con le tabelle e i
controlli pertinenti a quel particolare disegno sperimentale. \\
\addlinespace
\texttt{report\_doe.py} & Produce le statistiche descrittive e i confronti fra configurazioni in
forma di documento. \\
\addlinespace
\texttt{plot\_doe.py} & Genera le figure a partire dai risultati. \\
\addlinespace
\texttt{build\_relazione\_html.py} & Compone la versione consultabile della relazione sulle figure. \\
\bottomrule
\end{tabular}
\end{center}

\section{\texttt{1-configs/} --- le configurazioni di riferimento}

Contiene quattro configurazioni di carico in formato JSON, corrispondenti ai livelli di carico
impiegati negli esperimenti (\texttt{cfg\_n0}, \texttt{cfg\_n1}, \texttt{cfg\_n4},
\texttt{cfg\_n8}: il numero indica quanti task di disturbo affiancano il task critico), e un
documento che ne motiva le scelte, in particolare la modalit\`a di attivazione dei due tipi di task
e il valore della costante di calibrazione.

Queste configurazioni non vengono usate direttamente dalle campagne, che le rigenerano ogni volta a
partire dagli stessi parametri: servono a poter consultare e verificare il carico impiegato senza
dover eseguire nulla.

\section{\texttt{2-DoE/} --- dati sperimentali e risultati}

\`E la cartella pi\`u voluminosa e raccoglie l'intera evidenza sperimentale. Si articola in tre
livelli: le campagne, le celle, le singole esecuzioni.

\begin{center}\small
\begin{tabular}{p{3.5cm}ccp{7.6cm}}
\toprule
\textbf{campagna} & \textbf{celle} & \textbf{esecuzioni} & \textbf{fattore studiato} \\
\midrule
\texttt{block1/} & 4 & 80 & granularit\`a del tracciamento \\
\texttt{block2/} & 6 & 150 & criterio di campionamento \\
\texttt{block3/} & 12 & 180 & modalit\`a di consegna e carico concorrente \\
\texttt{diag/} & 3 & 75 & controllo di riproducibilit\`a \\
\bottomrule
\end{tabular}
\end{center}

\textbf{Le celle.} All'interno di ogni campagna, una cartella per combinazione di fattori, il cui
nome segue la stessa convenzione degli eseguibili con l'aggiunta del carico concorrente. Ad
esempio \texttt{t2\_p0\_s1\_r0.3\_e1\_n4} indica tracciamento di livello 2, consegna differita,
campionamento proporzionale al 30\,\%, uscita standard e quattro task di disturbo.

\textbf{Le esecuzioni.} Dentro ogni cella, una cartella \texttt{run\_NN} per ripetizione, che
contiene:

\begin{center}\small
\begin{tabular}{p{4.3cm}p{10.0cm}}
\toprule
\texttt{config.json} & La configurazione effettivamente eseguita, salvata accanto ai dati che ha
prodotto: ogni esecuzione \`e cos\`i autocontenuta e verificabile a posteriori. \\
\addlinespace
\texttt{rtapp-*.log.gz} & Un log per thread, con una riga per iterazione: durata della fase di
lavoro, periodo, istanti di inizio e fine, margine residuo rispetto alla scadenza e latenza di
risveglio. Sono la fonte primaria di tutte le misure. \\
\addlinespace
\texttt{stdout.log.gz} & Uscita standard; nelle esecuzioni configurate in tal senso contiene gli
elementi di traccia emessi, che vengono qui contati. \\
\addlinespace
\texttt{stderr.log.gz} & Uscita di errore, da cui si ricava fra l'altro il numero di tentativi di
trasmissione. \\
\bottomrule
\end{tabular}
\end{center}

I log sono conservati compressi: da circa 1.7\,GB complessivi si scende a poco pi\`u di 200\,MB,
mantenendo l'intera evidenza grezza a disposizione per verifiche successive.

\textbf{Risultati e sintesi}, al primo livello della cartella:

\begin{center}\small
\begin{tabular}{p{4.3cm}p{10.0cm}}
\toprule
\texttt{data\_table.csv} & Una riga per esecuzione con i livelli dei fattori e le condizioni
ambientali registrate (frequenza osservata, temperatura prima e dopo, esito). \`E il registro
prodotto durante le campagne. \\
\addlinespace
\texttt{results.csv} & Lo stesso registro arricchito con le variabili di risposta estratte dai log:
485 righe, una per esecuzione. \`E il file su cui si basa tutta l'analisi. \\
\addlinespace
\texttt{results\_summary.csv} & Una riga per cella, con le grandezze aggregate. \\
\addlinespace
\texttt{figures/} & Le figure prodotte dall'analisi. \\
\addlinespace
\texttt{REPORT.md}, \texttt{RELAZIONE\_GRAFICI.md} & Le sintesi discorsive dei risultati e la
lettura commentata di ciascuna figura. \\
\addlinespace
\texttt{NOTES.md} (in ogni campagna) & Le annotazioni specifiche di quella campagna: configurazione
impiegata, risultati e osservazioni. \\
\bottomrule
\end{tabular}
\end{center}

\section{Le cartelle rimanenti}

\textbf{\texttt{otel-installdir/}} contiene la libreria OpenTelemetry compilata e installata
localmente (sedici librerie statiche e le relative intestazioni), contro cui il generatore di carico
viene collegato. Non fa parte del codice del progetto e non viene versionata: va ricostruita dai
sorgenti della libreria se assente.

\textbf{\texttt{0-explore/}} raccoglie le prove preliminari eseguite manualmente prima di
automatizzare le misure, ciascuna con le proprie annotazioni. Sono servite a comprendere il
funzionamento del generatore e della strumentazione, a verificare l'efficacia dell'isolamento e a
individuare alcune insidie di misura poi tenute in conto nel disegno degli esperimenti. Non
contengono dati usati nelle conclusioni, ma documentano il percorso che vi ha condotto.

\textbf{\texttt{3-report/}} contiene i documenti di sintesi in formato PDF e i relativi sorgenti.

\vspace{.8em}
\hrule
\vspace{.6em}

\small
\textbf{In sintesi.} \texttt{rt-app/} e \texttt{otel-installdir/} sono il software sotto esame;
\texttt{bin/} ne conserva le versioni compilate, una per configurazione; \texttt{scripts/} contiene
gli strumenti che preparano la piattaforma, eseguono le campagne e ne analizzano i risultati;
\texttt{1-configs/} documenta i carichi impiegati; \texttt{2-DoE/} custodisce l'evidenza
sperimentale e le sue elaborazioni.

\end{document}
